\section{Method}
\label{sec:method}

In this section, we introduce the LV-MAE framework for learning long-video representations. Our approach operates hierarchically: first, we extract short-video representations using off-the-shelf models such as LanguageBind~\cite{zhu2024languagebind} and InternVideo2~\cite{Wang2024InternVideo2SV}. Next, we capture long-range dependencies across video segments through self-supervised learning. Specifically, we propose a masked-embedding autoencoder that operates on these short-video embeddings, with an architecture inspired by the asymmetric encoder-decoder design proposed in MAE~\cite{he2022masked}. In the following subsections, we delineate each component of the framework.

\subsection{Short-video Representation} \label{subsec:short_video_representation}
Given a full video $V$, we first segment it into a set of $N$ consecutive short videos (\eg, five seconds long), represented as $\mathcal{V} = \{\mathbf{v}_i\}_{i=1}^{N}$. For each segment, we extract $T$ frames, denoted by $\mathbf{v}_i \in \mathbb{R}^{T \times C \times H \times W}$, where $C$, $H$, and $W$ are the number of channels, height, and width of the frames, respectively. Each video segment $\mathbf{v}_i$ is then processed by a multimodal model, pre-trained for short video-text alignment, such as LanguageBind~\cite{zhu2024languagebind} or InternVideo2 \cite{Wang2024InternVideo2SV}. It produces a set of short-video embeddings $\mathcal{E} = \{\mathbf{e}_i\}_{i=1}^N$, where $\mathbf{e}_i \in \mathbb{R}^{d}$ and $d$ is the embedding dimension.

\subsection{Masked-embedding Autoencoder}

MAE~\cite{he2022masked} and its video-based variants~\cite{tong2022videomae, wang2023videomae} operate directly on video frames, masking a subset of image patches and reconstructing them at the decoder output. In contrast, our approach processes embedding vectors that represent sequential short video segments.

Formally, given a set of short-video embeddings $\mathcal{E}$, we mask a subset $\mathcal{M} \subset \mathcal{E}$, while only the remaining tokens $\mathcal{E} \setminus \mathcal{M}$ are passed as input to the encoder,
\begin{equation} \label{eq:encoding}
    \mathcal{Z} = \text{Encoder}(\mathcal{E} \setminus \mathcal{M}),
\end{equation}
where $\mathcal{Z}$ is the set of encoded visible embeddings. 
We then provide the decoder with both $\mathcal{Z}$ and $|\mathcal{M}|$ mask tokens, where each mask token is represented by a shared, learnable vector, denoted as $\mathbf{m}$. Additionally, we incorporate positional encoding at the encoder and decoder to indicate each token's location within the long video sequence. 

The decoder then outputs a set of reconstructed embeddings $\hat{\mathcal{E}} = \{\hat{\mathbf{e}}_i\}_{i=1}^N$, \ie, 
\begin{equation}
    \hat{\mathcal{E}} = \text{Decoder}\left(\mathcal{Z}, \mathbf{m}, \mathcal{P}\right),
\end{equation}
where $\mathcal{P}=\{\mathbf{p}_i\}_{i=1}^N$ is the set of positional embeddings added to both visible and mask tokens corresponding to their original locations.

\vspace{-4mm}
\paragraph{Loss function.}
We use mean squared error (MSE) loss during training, specifically computing the MSE between the original masked embeddings $\mathbf{e}_i \in \mathcal{M}$ and their corresponding reconstructed embeddings $\hat{\mathbf{e}}_i$ as
\begin{equation}
    \mathcal{L} = \frac{1}{|\mathcal{M}|} \sum_{\mathbf{e}_i \in \mathcal{M}} \left\lVert \mathbf{e}_i - \hat{\mathbf{e}}_i \right\rVert_2^2.
\end{equation}
Although simple, we find that the MSE loss is effective for pre-training long videos represented as sequences of embedding vectors.

\vspace{-4mm}
\paragraph{Masking.}
As previously mentioned, in LV-MAE, we mask a portion (\eg, 50\%) of the sequential embedding tokens by removing a subset $\mathcal{M}$ from the input set of tokens $\mathcal{E}$. In this work, we explore two masking strategies: \textit{random} masking and \textit{semantic} masking. In random masking, the subset $\mathcal{M}$ is simply generated by randomly sampling $|\mathcal{M}|$ embeddings and excluding them from $\mathcal{E}$.

Recently, several works have proposed alternative masking strategies based on masking salient regions within the spatio-temporal tokens~\cite{Li2022SemMAESM, Fan2023MotionGuidedMF, Huang2023MGMAEMG}. Inspired by these efforts, we propose a \textit{semantic} masking strategy that focuses on masking distinct or salient elements within the embedding sequence. Specifically, we employ a simple approach that leverages the semantic information within short-video representations. First, we compute the cosine similarity between each embedding $\mathbf{e}_i$ and its preceding embedding $\mathbf{e}_{i-1}$. Then, we select the embeddings with the lowest cosine similarity values and mask them. This encourages the model to learn more complex dependencies between video segments, enhancing its capacity for in-depth long-video understanding. We demonstrate that in some cases, the semantic strategy yields better results, particularly when employing linear probing for classification.


\vspace{-4mm}
\paragraph{Handling varying video lengths}
To handle videos of varying lengths and to maintain computational efficiency, we adopt techniques from BERT~\cite{devlin2019bert}. Each sequence of short-video segments is limited to a maximum of 256 tokens, with shorter sequences padded using a special \texttt{[PAD]} token. During the self-attention process, we employ an attention mask to prevent these padding tokens from influencing model training. 

\subsection{Training and Data} \label{subsec:train_and_data}
\paragraph{Diverse data sources.} Our model is pre-trained on diverse datasets across multiple video domains, enhancing its robustness and adaptability to various tasks. 
Specifically, we pre-train on over 1,000 long-length movies and TV series, encompassing a wide range of genres, styles, and lengths to ensure comprehensive content diversity.
In addition, we incorporate three publicly available datasets. 
\textit{FineVideo}~\cite{Farré2024FineVideo} includes a vast collection of diverse videos with more than 40,000 videos covering a wide spectrum of categories and domains. With thousands of hours of content, FineVideo represents a rich resource for capturing varied human activities, scenes, and interactions. Additionally, we leverage the train set of \textit{MovieClips}~\cite{wu2021towards} that contains approximately 7,000 video clips, typically ranging from one to three minutes, curated from various movies. It serves as a diverse collection of cinematic content. Finally, \textit{ActivityNet}~\cite{caba2015activitynet}, a dataset that focuses on complex human activities that are relevant to daily life, offering a broad range of action categories and scenarios. Our diverse video dataset supports training across varied domains, enhancing generalization. Since our framework is self-supervised, adding more data is easy and requires no manual labels.


\vspace{-4mm}
\paragraph{Training strategy.} To train our model, we employ a dynamic sampling strategy to handle videos of varying lengths. For each video, we randomly sample a duration in seconds, and this ensures that each batch contains videos of different lengths. It allows us to train on entire videos or portions of them, providing our model with exposure to both short and long video segments, which is particularly useful for handling complex temporal dependencies in long-form video understanding tasks. In practice, for simplicity, we set each short video segment to a length of five seconds. In App.~\ref{app:extended_lvu_ab} we explore the impact of segment length.

\vspace{-4mm}
\paragraph{Training efficiency.} Our pre-training method is highly efficient compared to existing approaches. This efficiency is demonstrated by the reduced number of tokens required during training. Specifically, since we leverage short-video representations extracted by off-the-shelf models, each video segment is represented by a single token. Thus, with each segment lasting five seconds, processing a two-minute video with LV-MAE requires only 24 tokens, compared to the 11,760 tokens required by other frame-level methods (60 frames $\times$ 14 $\times$ 14 patches for a standard 16 $\times$ 16 image grid)~\cite{li2024videomamba, tong2022videomae, wang2023videomae, islam2022long, lin2022learning}.
As a result, we can pre-train LV-MAE in only 20 hours on 8 $\times$ A10 GPUs while achieving state-of-the-art results on downstream tasks. 
Our optimized framework makes pre-training on long videos computationally feasible without compromising performance. Additional implementation details are provided in App.~\ref{app:implementation_details}.


\begin{table*}[h]
    \caption{\textbf{LVU benchmark results.} We compare the Top-1 accuracy results obtained by LV-MAE and other methods. ``FB'' (Frozen Backbone) indicates that only a small subset of parameters is fine-tuned, in contrast to other methods that tune the entire model. ``Dir.'' and ``Rel.'' refer to ``Director'' and ``Relationship,'' respectively. LanguageBind/InternVideo2-Baseline refers to using the raw embeddings.}
    \label{tab:lvu_benchmark}
    \centering
    \small
    \setlength{\tabcolsep}{6.5pt}
    \begin{tabular}[t]{l|c|cccc|ccc|c||cc}
            \hline 
            
            \hline
            
            \hline
        \multirow{2}*{\textbf{Method}} & \multirow{2}*{\shortstack{\textbf{FB}}} & \multicolumn{4}{c|}{\textbf{Metadata $\uparrow$}} & \multicolumn{3}{c|}{\textbf{Content $\uparrow$}} & \multirow{2}*{\textbf{Avg.}} & \multicolumn{2}{c}{\textbf{User $\downarrow$}} \\
        & & \textbf{Dir.} & \textbf{Genre} & \textbf{Writer} & \textbf{Year} & \textbf{Scene} & \textbf{Speak} & \textbf{Rel.} & & \textbf{Likes} & \textbf{Views} \\
        \hline
        VideoBERT~\cite{sun2019videobert} & \xmark & 47.30 & 51.90 & 38.50 & 36.10 & 54.90 & 37.90 & 52.80 & 45.6 & 0.32 & 4.46 \\
        Object Transformer~\cite{wu2021towards} & \xmark & 51.20 & 54.60 & 34.50 & 39.10 & 56.90 & 39.40 & 53.10 & 46.9 & \underline{0.23} & 3.55 \\
        LST~\cite{islam2022long} & \xmark & 56.07 & 52.70 & 42.26 & 39.16 & 62.79 & 37.31 & 52.38 & 48.9 & 0.31 & 3.83 \\
        Performer~\cite{islam2022long} & \xmark & 58.87 & 49.45 & 48.21 & 41.25 & 60.46 & 38.80 & 50.00 & 49.6 & 0.31 & 3.93 \\
        Orthoformer~\cite{islam2022long} & \xmark & 55.14 & 55.79 & 47.02 & 43.35 & 66.27 & 39.30 & 50.00 & 50.9 & 0.29 & 3.86 \\
        ViS4mer~\cite{islam2022long} & \xmark & 62.61 & 54.71 & 48.80 & 44.75 & 67.44 & \underline{40.79} & 57.14 & 53.7 & 0.26 & 3.63 \\
        VideoMamba~\cite{li2024videomamba} & \xmark & \underline{67.29} & 65.24 & 52.98 & 48.23 & 70.37 & 40.43 & \textbf{62.50} & 58.1 & 0.26 & 2.90 \\
        \hline
        LanguageBind-Transformer & \xmark & 24.30 & 57.88 & 16.07 & 18.44 & 35.80 & 32.45 & 51.22 & 33.7 & 0.56 & 3.04 \\
        LanguageBind-Mamba & \xmark & 61.68 & \underline{70.38} & 51.78 & \underline{56.74} & \underline{74.04} & 38.83 & 43.90 & 56.8 & 0.30 & 2.98 \\
        InternVideo2-Baseline~\cite{Wang2024InternVideo2SV} & \cmark & 45.79 & 54.79 & 5.36 & 8.51 & 34.57 & 30.32 & 51.22 & 32.9 & 0.24 & 3.22 \\
        LanguageBind-Baseline~\cite{zhu2024languagebind} & \cmark & 64.48 & 61.47 & 5.95 & 29.08 & 30.86 & 31.91 & 48.78 & 38.9 & \textbf{0.21} & 2.90 \\
        \hline
        \rowcolor{Gray!15} $\text{LV-MAE}_{\text{InternVideo2}}$(Ours) & \cmark & 62.62 &  68.15 & \underline{57.14} & \textbf{58.15} & \textbf{77.78} & 39.89 & 53.66 & \underline{59.6} & \underline{0.23} & \underline{2.68} \\
        \rowcolor{Gray!15} $\text{LV-MAE}_{\text{LanguageBind}}$(Ours) & \cmark & \textbf{77.57} & \textbf{71.57} & \textbf{64.28} & \textbf{58.15} & 72.84 & \textbf{40.95} & \underline{58.53} & \textbf{63.4} & \underline{0.23} & \textbf{2.52} \\
            \hline 
            
            \hline
            
            \hline
    \end{tabular}
    \vskip -0.1in
\end{table*}



\begin{table}[h]
    \caption{\textbf{Breakfast and COIN benchmark results}. We compare the Top-1 accuracy results obtained by LV-MAE and other methods.  ``FB'' (Frozen Backbone) indicates that only a small subset of parameters is fine-tuned, in contrast to other methods that tune the entire model. Baseline refers to using the raw embeddings.}
    \vspace{-1.5em}
    \begin{center}
    \small
        \begin{tabular}{l|c|cc}
            \hline
            
            \hline
            
            \hline\\[-3mm]
            {\multirow{2}{*}{\textbf{Method}}} & {\multirow{2}{*}{\textbf{FB}}} & \textbf{Breakfast} & \textbf{COIN}\\
            & & Top-1 & Top-1 \\
            \hline
            Timeception~\cite{hussein2019timeception} & \xmark & 71.3 & - \\
            VideoGraph~\cite{hussein2019videograph} & \xmark & 69.5 & - \\
            GHRM~\cite{zhou2021graph} & \xmark & 75.5 & - \\
            Distant Supervision~\cite{lin2022learning} & \xmark & 89.9 & 90.0 \\
            ViS4mer~\cite{islam2022long} & \xmark & 88.2 & 88.4 \\
            Turbo~\cite{Han22b} & \xmark & 91.3 & 87.5 \\
            VideoMamba~\cite{li2024videomamba} & \xmark & \textbf{96.9} & 90.4 \\

            \hline
            InternVideo2-Baseline~\cite{Wang2024InternVideo2SV} & \cmark & 83.94 & 88.41 \\
            LanguageBind-Baseline~\cite{zhu2024languagebind} & \cmark & 73.52 & 91.13 \\
            \hline
            \rowcolor{Gray!15} $\text{LV-MAE}_{\text{InternVideo2}}$ (Ours) & \cmark & \underline{93.24} & \textbf{92.72} \\
            \rowcolor{Gray!15} $\text{LV-MAE}_{\text{LanguageBind}}$ (Ours) & \cmark & 91.55 & \underline{92.42} \\
            \hline 
            
            \hline
            
            \hline
        \end{tabular}
    % }
    \end{center}
\vspace{-1.2em}
\label{tab:bf_coin_benchmark}
\vspace{-0.5em}
\end{table}




\subsection{Interpretable Predictions} \label{subsec:interpretability}

One potential limitation in training masked-embedding autoencoders, where we process sequences of embeddings rather than spatio-temporal image patches at the frame level, is that it can be challenging to directly interpret the quality of token reconstruction produced by the decoder (\eg as done in \cite{he2022masked}). To address this, we leverage the aligned subspace of language-video embeddings from the pre-trained multimodal model to interpret predictions and assess the model's ability to reconstruct embeddings with accurate semantic meaning.
Specifically, we propose using retrieval for each reconstructed masked token against a large set of captions to visualize and evaluate the matches.
\vspace{-5mm}
\paragraph{Constructing caption database.} We begin by annotating a large set of short video segments from the LVU dataset, utilizing Claude-3.5 Sonnet \cite{anthropic2024claude35sonnet} to automatically generate textual descriptions for each segment. Specifically, we uniformly sample 20 frames from each segment and prompt the LLM to provide a concise description that captures the key visual or semantic information in the segment. Formally, for each segment $v_i$, we obtain an annotation $a_i$.
Next, we use LanguageBind's language encoder to extract textual embedding for each segment's description, $\mathbf{e}_i^{\ell}$.
\vspace{-5mm}
\paragraph{Masked-embeddings retrieval.} Given a reconstructed embeddings output by the decoder $\hat{\mathbf{e}}_j$, corresponding to the $j$-th video segment, we retrieve the top-matching captions with the highest response by computing the cosine similarity between $\hat{\mathbf{e}}_j$ and each caption's textual embedding,
$\mathbf{e}_i^{\ell}$. This way, these top-matched retrieved captions can be inspected to ensure that the semantic meaning of the reconstructed masked embeddings is well captured by the decoder. Visualization examples are provided in section \ref{sec:exp_interpretabitlity}.
